\scanpage{359}
taking traces we find
\[
  2\lambda=2 \tag{45}
\]

--- these are the necessary and sufficient conditions on $\lambda$ for
\deleted{\unsure{$\lambda$}}
($\varepsilon$ being regular unipotent)
$\lambda\varepsilon$ to be still \inserted{\unsure{strict\textsuperscript{ly}}}
\unsure{unipoten\textsuperscript{t}} (hence necessari\textsuperscript{ly} regular unipotent).
If $2$ is invertible on $S$, \inserted{\uline{more generally}
if $2$ is a regular section of $\underline{\mathcal O}_S$,}
\deleted{\unsure{is reduced}}
these conditions imply $\lambda=1$, OK, but
\deleted{\unsure{in the case}}\inserted{outside} of this hypothesis (for example in characteristic $2$
\deleted{\unsure{even if}} over a nonreduced scheme) this is no longer necessarily so. Let us write
\[
  \lambda=1+\nu
\]
the conditions (44) (45) read
\[
  2\nu=0,\quad \nu^2=0, \tag{46}
\]
thus for the map (43) to be an injection on $S$, it is necessary and sufficient
that the \inserted{only} sections of $S$ which satisfy (46)
are the zero section.

In general, denoting by $\mu'_2$ the subgroup scheme
of $\mu_2$ defined by (45)
(it is in fact an absolute infinitesimal group \inserted{scheme},
\deleted{\unsure{over}} \inserted{finite over $\mathbb Z$, not flat over $\mathbb Z$,}
isomorphic to the spectrum of
\[
  \mathbb Z[T]/(T^2-1,2T-2)\simeq
  \mathbb Z[S]/(S^2,2\unsure{S}),
\]
also isomorphic to the scheme of dual numbers over $\mathbb Z$,
relative to the module $\mathbb Z/2\mathbb Z$ over $\mathbb Z$,
\deleted{\unsure{$\mu'_2\simeq$}}
\[
  \mu'_2\simeq\operatorname{Spec}
  \bigl(\mathbb Z\deleted{\unsure{\varepsilon}}\oplus\underbrace{\mathbb Z/2\mathbb Z}_{\text{square-zero ideal}}\bigr)
  \quad\ldots
\]
\mspage{358}

% Printed archive footer outside the manuscript frame:
% 358
% https://grothendieck.umontpellier.fr
% Université de Montpellier, UMR CNRS 5149, Case courrier 051, Place Eugène Bataillon
% 34095 Montpellier cedex 5 - France

\scanpage{360}
\mspage{481}
It turns out that $\mu'_{2S}$ operates
\inserted{freely by translations}
on the scheme of regular unipotents \inserted{\unsure{(or not)}}
of $G$, and we have
\[
\diagram{\begin{tikzcd}[ampersand replacement=\&]
\underline{\mathrm{Unip.rég}}(G)/\mu'_{2S}
  \arrow[r,"\sim"] \& \underline{\mathrm{Unip.rég}}(PG)
\end{tikzcd}}\ . \tag{47}
\]
This seems at first sight to be in contradiction with the \unsure{sequence} of Borels
\[
\diagram{\begin{tikzcd}[ampersand replacement=\&,column sep=small,row sep=large]
\deleted{\unsure{1}} \& B \arrow[d,mapsto]
 \& \underline{\mathrm{Borels}}(G)
   \arrow[r,"\sim","B\mapsto B/\mathbb \mathbb G_m"']
   \arrow[d,"\wr"']
 \& \mathrm{Borels}(PG) \arrow[d,"\wr"]
 \& \deleted{\unsure{B}}\widetilde B \arrow[d,mapsto] \\
 \& B_u
 \& \begin{gathered}\text{``Subgroups}\\
     \deleted{\unsure{\mathrm U}}\text{unipotent with 1}\\
     \text{parameter'' of }G\end{gathered}
     \arrow[r,"\sim"]
 \& \begin{gathered}\text{``Subgroups}\\
     \text{unipotent with 1 parameter''}\\\text{of }PG\end{gathered}
 \& \widetilde B_u \\
 \& \& \text{\unsure{that is}}\quad L \arrow[r,mapsto,"\sim"] \& p(L)
\end{tikzcd}} \tag{48}
\]
(where $p:G\to PG$ is the canonic\textsuperscript{al} morphism)

Could the difficulty come from the fact that the vertical arrows
(surjective, by definition of the notion of ``1-parameter unipotent subgroups''
of $G$ or of $PG$) might not be injective, outside the hypothesis that $2$ is invertible
(i.e. \unsure{residue} characteristic\ $\ne2$)?
I feel like saying that
\[
  \underline B=\operatorname{Norm}(B_u),\quad
  \widetilde B=\operatorname{Norm}(\widetilde B_u) \tag{49}
\]
are formulas which allow us to recover $B,\widetilde B$ with the help
of $B_u,\widetilde B_u$\unsure{,|''} --- and I have
\deleted{\unsure{done}} indeed \uline{\unsure{already}}
checked the first of these formulas, by dint of linear algebra in $M$
(in the form: the 1-parameter unipotent subgroups of $G$
correspond $1$-$1$ to the sections of $\mathbb P$ over $S$,
i.e. to the line \unsure{subbundles} $D$ of $F$\ldots) ---
\mspage{359}

% Printed archive footer outside the manuscript frame:
% 359
% https://grothendieck.umontpellier.fr
% Université de Montpellier, UMR CNRS 5149, Case courrier 051, Place Eugène Bataillon
% 34095 Montpellier cedex 5 - France

\scanpage{361}
but \unsure{this same} point in $PG$, it is
\deleted{if $u\in\Gamma M^*$ is such that
$u\widetilde N u^{-1}=\mu N+\lambda$,}
\inserted{not very trivial.}
\deleted{with $\mu\in\Gamma\mathbb \mathbb G_m$,}

The difficulty must rather come from the fact \unsure{that a|that no}
\unsure{regular unipotent} section $\widetilde\varepsilon$ of $PG$
does not define (without restriction on $S$) a 1-parameter unipotent subgroup
--- whereas for a section $\varepsilon$ of
\unsure{regular unipotents|regular unipotent} of $G$
\[
  \varepsilon=1+N,
\]
we define \inserted{one} by the formula
\[
  \varepsilon(t)=1+tN\ .
\]
We \unsure{will note} that if we replace $\varepsilon$ by
$\inserted{\varepsilon'=}\lambda\varepsilon$
(with $\lambda$ satisfying (44) (45) $\lambda^2=1$, $2\lambda=2$),
We have
\[
  \varepsilon'=(1+\nu)(1+N)=1+N'
  \qquad N'=(\unsure{1+\nu})N+\nu
\]
and
\[
  \varepsilon'(t)=1+tN'
\]
I say that if \inserted{$\lambda\ne1$ i.e.}
$\nu\ne0$, \deleted{\unsure{we then have}} the composite homomorphisms
\[
\diagram{\begin{tikzcd}[ampersand replacement=\&]
\widetilde{\underline\varepsilon},\widetilde{\underline\varepsilon}'\colon
\mathbb G_{\unsure{a|0}} \arrow[r] \& G \arrow[r] \& PG
\end{tikzcd}}
\]
\margin{$\Vert$}
\deleted{\unsure{will be}} associated to $\varepsilon,\varepsilon'$
will be \uline{distinct}, although they have
the same value for $t=1$.
Indeed, \deleted{\unsure{we will have}} the derivatives at the origin
of $\widetilde{\underline\varepsilon},\widetilde{\underline\varepsilon}'$
are identified with the composites
\[
\diagram{\begin{tikzcd}[ampersand replacement=\&]
\underline{\mathcal O}_S \arrow[r] \& M \arrow[r] \& M/\underline{\mathcal O}_S\unsure{=|\simeq}P\unsure{\mathfrak g}
\end{tikzcd}}
\]
\mspage{360}

% Printed archive footer outside the manuscript frame:
% 360
% https://grothendieck.umontpellier.fr
% Université de Montpellier, UMR CNRS 5149, Case courrier 051, Place Eugène Bataillon
% 34095 Montpellier cedex 5 - France

\scanpage{362}
\mspage{482}
which send $1$ to $\widetilde N$ resp. to $\widetilde N'$,
(where the $\sim$ here denotes passage to the \uline{additive} quotient) we have
\[
  \widetilde N'=(1+\nu)\widetilde N
\]
and $\widetilde N$ is a section of $P\unsure{\mathfrak g}$ everywhere nonzero
--- thus we can have \deleted{not}
\[
  \widetilde N'=\widetilde N\quad\text{i.e.}\quad
  (1+\nu)\widetilde N=\widetilde N\quad\text{i.e.}\quad
  \nu\widetilde N=0
\]
only if $\nu=0$, i.e. $\lambda=1+\nu=1$, q.e.d.

In conclusion, we must consider that we give ourselves in $PG$,
not only $\widetilde\rho$ and $\widetilde\sigma$
(whence $\widetilde\varepsilon_0=\widetilde\sigma^{-1}\widetilde\rho$,
$\widetilde\varepsilon_1=\widetilde\rho\widetilde\sigma^{-1}$),
but also ``1-parameter subgroups'' (parametrized)
\[
\diagram{\begin{tikzcd}[ampersand replacement=\&]
\deleted{\mathbb G_{\unsure{a|0}}}
  \arrow[r,"\deleted{\unsure{\widetilde{\underline\varepsilon}_0}}"] \& \deleted{PG}
\end{tikzcd}}\ ,
\]
\[
\left\{\diagram{\begin{tikzcd}[ampersand replacement=\&,row sep=small]
\widetilde{\underline\varepsilon}_0\colon \mathbb G_{\unsure{a|0}}
  \arrow[r] \& PG \\
\widetilde{\underline\varepsilon}_1\colon \mathbb G_{\unsure{a|0}}
  \arrow[r] \& PG
\end{tikzcd}}\right.\ . \tag{49}
\]
related of course by
\[
\left\{\begin{aligned}
\widetilde{\underline\varepsilon}_1
 &=\rho(\widetilde{\underline\varepsilon}_0)
   \bigl(=\sigma(\widetilde{\underline\varepsilon}_0)\bigr)
   \quad\text{i.e.}\\
\widetilde{\underline\varepsilon}_1(t)
 &=\rho\bigl(\widetilde{\underline\varepsilon}_0(t)\bigr)
   =\sigma\bigl(\widetilde{\underline\varepsilon}_0(\unsure{t|u})\bigr)
\end{aligned}\right. \tag{50}
\]
which means precisely that we \uline{give} ourselves, in addition to
$\widetilde\rho,\widetilde\sigma$, \unsure{a} lifting\deleted{\illegible{1}}
\inserted{\inserted{regular} \unsure{unipoten\textsuperscript{t}}}
$\varepsilon_0$ of $\widetilde\varepsilon_0$
(\deleted{\unsure{which|that}} the choice is therefore not \unsure{a} priori unique),
and we deduce $\varepsilon_1$ from it by
\[
  \deleted{\unsure{\rho}}\quad
  \varepsilon_1=\widetilde\rho(\varepsilon_0)
  =\widetilde\sigma(\varepsilon_0)
\]
\margin{we will necessarily have a relation
$\sigma^{-1}\rho= -\lambda_0\varepsilon_0$
hence \deleted{$\unsure{\operatorname{int}(\varepsilon_0)|\mathrm{i.e.}\,\sigma(\varepsilon_0)}\,\unsure{\varepsilon_0}$}
$\unsure{\sigma(\varepsilon_0)|\sigma^{-1}(\varepsilon_0)}=\rho(\varepsilon_0)\,\unsure{=}$}
\mspage{361}

% Printed archive footer outside the manuscript frame:
% 361
% https://grothendieck.umontpellier.fr
% Université de Montpellier, UMR CNRS 5149, Case courrier 051, Place Eugène Bataillon
% 34095 Montpellier cedex 5 - France

\scanpage{363}
\mspage{4}
\deleted{\unsure{And}}
Thus, we have adjusted the \deleted{\unsure{precision|position}}
\inserted{description} of the data, so as to
\deleted{\unsure{adjust}} be able to give precise normalization
statements --- namely

\uline{Proposition} Let $\widetilde\rho,\widetilde\sigma$
be sections of $PG$, and $\widetilde{\underline\varepsilon}_0$
a ``strict'' \inserted{additive} subgroup \unsure{with} 1 parameter
\margin{(i.e. a monomorphism and whose image is a $\widetilde B_u$)}
\[
\diagram{\begin{tikzcd}[ampersand replacement=\&]
\widetilde{\underline\varepsilon}_0\colon
\mathbb G_{\unsure{a|0}}
  \arrow[r,phantom,"\unsure{\longrightarrow|\text{---}}" description] \& PG
\end{tikzcd}}\ .
\]
\unsure{such|such} that we have \unsure{$\times$|$*$}
\[
  \widetilde\sigma^{-1}\widetilde\rho
   =\widetilde{\underline\varepsilon}_0(1)
   \overset{\mathrm{dfn}}=\widetilde\varepsilon_0 \tag{51}
\]
whence
\[
\left\{\begin{aligned}
\widetilde\sigma(\widetilde{\underline\varepsilon}_0)
 &=\widetilde\rho(\widetilde{\underline\varepsilon}_0)
   \overset{\mathrm{dfn}}=\widetilde{\underline\varepsilon}_1,
   \quad\text{and}\\
\widetilde\rho\widetilde\sigma^{-1}
 &=\widetilde{\underline\varepsilon}_1(1)
   \overset{\mathrm{dfn}}=\widetilde\varepsilon_1
   =\rho(\widetilde\varepsilon_0)=\sigma(\widetilde\varepsilon_0)
\end{aligned}\right. \tag{52}
\]

Then:

1) (\unsure{``For the record''})
$\exists$ unique $\varepsilon_0,\varepsilon_1$
\inserted{\deleted{\unsure{strict\textsuperscript{ly}}} regular unipotent} sections
\deleted{such} of $G$ such that
(setting $\varepsilon_i=\deleted{\unsure{\mathrm{Id}+}}1+N_i$,
$\underline\varepsilon_i(t)=1+tN_i$) we have
\[
  \widetilde{\underline\varepsilon}_i(t)=\varepsilon_i(t)^{\sim}\ . \tag{53}
\]

2) \deleted{Suppose that $\varepsilon_0$}
\inserted{We will then have}
\[
  \unsure{\dot{\varepsilon}_1|\varepsilon_1}=\widetilde\sigma(\varepsilon_0)
    =\widetilde\rho(\varepsilon_0) \tag{54}
\]
\[
\deleted{
  \unsure{\rho}\sigma^{-1}\rho=-\lambda_0\varepsilon_0
  \quad\varepsilon_0=-\lambda_0\sigma^{-1}\rho,\quad
  \varepsilon_1=-\lambda_0\rho\sigma^{-1},
}\tag{\deleted{55}}
\]
\deleted{3) For $\varepsilon_0$ and $\varepsilon_1$ to be
everywhere independent \unsure{in} each fiber, i.e. $N$}

3) \unsure{Let|Being} $\rho,\sigma$ be liftings of
$\widetilde\rho,\widetilde\sigma$, we will then have
\[
  \deleted{\unsure{\rho\sigma^{-1}\rho}}\qquad
  \varepsilon_0=-\lambda_0\sigma^{-1}\rho,\quad
  \varepsilon_1=-\lambda_0\rho\sigma^{-1} \tag{55}
\]
for $\lambda_0\in\Gamma\mathbb \mathbb G_m$ uniquely determined.
For a suitable choice of $\rho$ (when $\sigma$ is already chosen)
we can assume $\lambda_0=1$
(which determines $\rho$ uniquely in
\deleted{\unsure{function}} terms of $\sigma$, \unsure{and} conversely)
\mspage{362}

% Printed archive footer outside the manuscript frame:
% 362
% https://grothendieck.umontpellier.fr
% Université de Montpellier, UMR CNRS 5149, Case courrier 051, Place Eugène Bataillon
% 34095 Montpellier cedex 5 - France